\section{Step 91 Admit picture and macroblock controls}\label{h263-syntax-s01}
Validate records, complete extents and full-width selectors before level reads.
Reject INTER in an I-picture. Reject nonzero vector/predictor fields for INTRA;
they have no syntax here and usually indicate a mismatched request. For INTER,
validate all four signed half-sample components before computing differences.

\paragraph{Observe} An INTRA macroblock in a P-picture is valid, but it uses the
P-picture INTRA MCBPC group and has COD zero before it.
\Needspace{8\baselineskip}
\paragraph{Check} Which picture/type combination is excluded?
\begin{itemize}
  \item[A.] INTER inside an I-picture.
  \item[B.] INTRA inside a P-picture.
  \item[C.] INTRA inside an I-picture.
\end{itemize}

\section{Step 92 Validate all six blocks and form presence}\label{h263-syntax-s02}
Scan all 384 halfwords before invoking output children. Validate each block's
special DC entry for INTRA and the ordinary level range elsewhere. Set the
appropriate presence bit for each nonzero ordinary value, excluding INTRA DC.
Initialize the presence accumulator afresh on every call.

\paragraph{Observe} Six INTRA DC codes 255 with no AC values give mask zero.
A forbidden final coefficient in block six still returns status one before output.
\Needspace{8\baselineskip}
\paragraph{Check} Does a nonzero INTRA DC code by itself set that block's flag?
\begin{itemize}
  \item[A.] Yes, every DC value sets every chroma flag.
  \item[B.] No, this presence decision concerns ordinary AC levels for INTRA.
  \item[C.] Only when the host uses little endian.
\end{itemize}

\section{Step 93 Distinguish skipped from coded zero motion}\label{h263-syntax-s03}
A skip candidate requires a P-picture, INTER mode, zero presence mask, selected
motion zero and permission to skip. The predictor need not be zero. An eligible
skip emits COD one and nothing else. If skip is disallowed or not chosen, build
ordinary coded INTER syntax and encode the difference needed to recover the
selected vector, even when that selected vector is zero.

\paragraph{Observe} Predictor 31 and selected vector zero require a nonzero MVD
in coded output, while COD one directly declares actual zero motion.
\Needspace{8\baselineskip}
\paragraph{Check} What motion does a baseline skipped macroblock use?
\begin{itemize}
  \item[A.] The supplied predictor without modification.
  \item[B.] The preceding raw picture's arbitrary vector.
  \item[C.] Zero.
\end{itemize}

\section{Step 94 Select MCBPC and CBPY tables}\label{h263-syntax-s04}
Select the MCBPC group from picture type and INTRA/INTER mode, then index its
four entries with the low two presence bits. Form the direct Y pattern from
the high four bits. For INTRA, use that direct CBPY table index. For INTER, use
its four-bit complement. Retain both code values and widths in the private header plan.

\paragraph{Observe} Mask 33 describes Y block one and Cr block six. CBPC is one,
and the direct Y pattern is eight; the INTER CBPY lookup index is seven.
\Needspace{8\baselineskip}
\paragraph{Check} Which lookup index represents an INTER Y pattern of zero?
\begin{itemize}
  \item[A.] Fifteen.
  \item[B.] Always zero.
  \item[C.] The chroma stride.
\end{itemize}

\section{Step 95 Normalize both motion differences}\label{h263-syntax-s05}
For each component, subtract predictor from selected half-sample motion, add
32, retain the low six bits and use that value as the table index. The index
represents normalized difference plus 32. Emit horizontal before vertical.
Keep the original signed components for validation rather than truncating their
high bits before checking the admitted range.

\paragraph{Observe} Selected minus 32 and predictor 31 give raw difference minus
63, normalized difference one, and decoder recovery of minus 32 after wrapping.
\Needspace{8\baselineskip}
\paragraph{Check} What are the vector units at this boundary?
\begin{itemize}
  \item[A.] Always full pixels regardless of the earlier search.
  \item[B.] Half-sample units; full-pixel search results are multiplied by two.
  \item[C.] Bytes in the output file.
\end{itemize}

\section{Step 96 Emit the header to private storage}\label{h263-syntax-s06}
Initialize a zeroed private byte buffer and writer context at the caller's bit
offset modulo eight. Emit the planned COD/type/presence/motion fields, checking
every child return. Skip mode stops after COD one. Coded mode continues with
exactly the fields required by the selected type.

A zeroed staging buffer makes every new bit deterministic without reading the
caller yet. Retain enough space for the complete worst-case supported macroblock,
not just the header.

\paragraph{Observe} The default I-INTRA DC-only header has five bits before the
six mandatory DC bytes. It has no COD or MVD fields.
\Needspace{8\baselineskip}
\paragraph{Check} Which fields are absent from an I-INTRA header here?
\begin{itemize}
  \item[A.] MCBPC and CBPY.
  \item[B.] All six DC bytes.
  \item[C.] COD and motion-vector differences.
\end{itemize}

\section{Step 97 Emit block data in component order}\label{h263-syntax-s07}
Visit the four Y blocks, Cb and Cr in the published order. INTRA invokes the
coefficient emitter for every block, including those without AC. INTER invokes
it only for present blocks. Pass the same private context across these calls so
payloads continue at the actual next bit without automatic byte alignment.

\paragraph{Observe} INTRA presence zero still emits forty-eight DC bits. INTER
presence zero emits no coefficient blocks, though a coded header still has MVD.
\Needspace{8\baselineskip}
\paragraph{Check} Why must an INTRA block without AC still be emitted?
\begin{itemize}
  \item[A.] Its separate DC byte remains mandatory.
  \item[B.] It must contain an invented end-of-block code.
  \item[C.] Every block is padded to a native doubleword.
\end{itemize}

\section{Step 98 Check total space and merge boundary bits}\label{h263-syntax-s08}
After every child succeeds, subtract the initial private bit offset from the
private context count to obtain appended length. Check the caller's remaining
space before any caller-buffer load. Merge existing prefix bits into the private
first byte and unused suffix bits into its last byte. When the touched window
has one byte, both masks apply to that same staged byte.

Publish only touched bytes and commit the new caller count afterward. A failure
before that point returns count zero and retains all caller state.

\paragraph{Observe} A one-bit skipped append preserves seven other byte bits,
including the case where it ends exactly at the byte boundary.
\Needspace{8\baselineskip}
\paragraph{Check} When may the routine first read caller output bytes?
\begin{itemize}
  \item[A.] Before validating the last block.
  \item[B.] After all children succeed and the complete representation fits.
  \item[C.] After checking only MCBPC length.
\end{itemize}

\section{Step 99 Decode complete macroblock meaning}\label{h263-syntax-s09}
Run public feedback on a source with completed dependencies:
\begin{verbatim}
python3 construction/h263-macroblock-syntax/grade.py \
    my-macroblock.asm --cli /path/to/mars-cli.jar \
    --report macroblock-results.json
\end{verbatim}
The independent decoder checks type, presence mask, selected motion, all six
level arrays and exact payload consumption. It accepts legal compact/escape
choices and permitted skipped or coded zero-motion output. State and memory traces
check complete preflight, rejection, boundary preservation and saved registers
in both native byte orders.

\paragraph{Observe} Passing one header value is insufficient: a correct header
with omitted INTRA DC or swapped Cb/Cr payloads still fails decoded meaning.
\Needspace{8\baselineskip}
\paragraph{Check} Which result defines semantic agreement here?
\begin{itemize}
  \item[A.] The reference's exact private bytes.
  \item[B.] Only the printed number of bits.
  \item[C.] Decoded type, motion, flags and all selected block levels.
\end{itemize}

\section{Step 100 Leave picture history to the picture controller}\label{h263-syntax-s10}
Exercise repeated macroblock calls with the retained writer context, changing
type, levels and predictor fields between calls. Check that presence and header
plans are fresh. The picture controller must compute neighbour predictors, place
reconstructed blocks in Y/Cb/Cr planes and retain actual motion for future neighbours.
That history uses zero for INTRA and skipped macroblocks under this baseline path.

This emitter changes no picture or motion-history array. Picture-level ownership,
GOB boundaries, start-code alignment, file output and independent full-stream
checking are still separate integrations. Keep these boundaries explicit when
measuring the complete encoder rather than a syntax fragment.

\paragraph{Observe} A permitted skip and an equivalent coded zero-motion macroblock
can have different lengths while both declare actual motion zero and zero residuals.
\Needspace{8\baselineskip}
\paragraph{Check} Who supplies the neighbour-derived motion predictor?
\begin{itemize}
  \item[A.] The picture controller, using the encoded macroblock history.
  \item[B.] Unused suffix bits in the output buffer.
  \item[C.] The quantizer's DC code alone.
\end{itemize}
